\section{Related work}\label{sec:related}

\subsection{Leakage as a reproducible failure mode}

Kapoor and Narayanan~\cite{kapoor2023leakage} catalogue leakage across machine-learning-based science and give it
a working taxonomy: leakage occurs when information that would not be available at prediction
time enters training, and it is dangerous precisely because the result looks better rather than
worse. One of their categories, ``non-independence between train and test'', covers our case
directly: two rows derived from the same underlying unit are not independent, and a random
split over rows therefore delivers an optimistic estimate of generalisation. The same
diagnosis has been made in recommendation \citep{ji2023leakage}, where offline evaluation
routinely mixes user interactions across the split, and in translational
EEG \citep{brookshire2024eeg}, where segments of one recording are treated as separate samples. The phenomenon also has a data-mining formulation: leakage arises when a feature is generated
after the target and therefore carries information about it \citep{kaufman2012}. Our case is a
different mechanism with the same consequence. The window is not generated after the label; it
comes from the same recording, so a training row and a test row can be two views of one
performance and are not independent draws.

The methodological implication is that the unit of splitting must be named and, when the release
format makes rows rather than units look atomic, verified against the files. A benchmark that
releases feature images falls into exactly this trap, and the trap is silent: the pipeline runs
without error, the score looks plausible, and only the split unit distinguishes an optimistic
number from a transferable one.

\subsection{Splitting practice in audio}

Recording-level separation is an established practice in audio evaluation. Slice- or chunk-level
partitioning can expose a model to siblings of test examples, making within-recording recognition
possible rather than ensuring transfer to new recordings. Salamon et al.~\cite{salamon2014urbansound}
formalised this for urban sound, fixing ten folds at the recording level so that folds are
comparable across papers. Corpora that share the practice differ in what they document as the
splitting unit rather than in a single convention: FMA fixes an 8:1:1 partition at the recording
level \citep{defferrard2016}, whereas tagging and transfer studies built on curated track-level
collections report track identifiers instead \citep{choi2017,lee2017,won2020}, and shared
evaluation campaigns have made cross-paper comparison routine \citep{downie2010mirex}; CCMusic
extends this landscape to a multi-library Chinese-music resource \citep{zhou2025ccmusic}.
GiantMIDI-Piano is instead a large classical-piano MIDI resource \citep{kong2022giantmidi}.
Environmental-sound datasets and augmentation studies provide useful evaluation recipes
\citep{piczak2015,salamon2017augment}; we do not take them as evidence that augmentation
quantitatively removes, or fails to remove, leakage. The
difficulty is sharpest for collections distributed as pre-computed features: the audio may or may
not accompany the images, the file names encode windows, and new users inherit whatever split the
release implies.

\subsection{Released representations and what they cost}

Release formats changed as audio models grew. Large annotated corpora of audio events and music \citep{gemmeke2017audioset,engel2017} made it attractive to publish embeddings or spectrograms instead of waveforms \citep{hershey2017cnn,kong2020panns}, and self-supervised audio models now ship representations that downstream users fine-tune rather than retrain \citep{elizalde2024clap,li2023}. Each of these releases moves the split decision upstream: whoever pre-computes the features also decides what an item is, and a downstream user who sees one image per row has no reason to suspect that rows are windows.

\subsection{Dataset audits}

Auditing a published dataset is an established genre. Sturm~\cite{sturm2012gtzan} showed that a
widely used genre benchmark contains mislabels and duplicates that change conclusions drawn
from it. Our audit differs in object: we do not question the content of the labels or the
duplication structure of the corpus, but the protocol that maps that corpus onto train and test.
Audits of this kind have a second, sharper form: a system can be shown not to listen at all, in which case the benchmark measures the shortcut rather than the task \citep{sturm2014horse}. Re-measuring a leaderboard number under a different protocol and finding that it moves is likewise an established result \citep{recht2019}, and dataset documentation now asks authors to state the split unit rather than leave it to the reader \citep{gebru2018}. That distinction matters for how the result should be read --- a labelling defect is a property
of the data, whereas a split-unit defect is a property of the evaluation recipe and can be
repaired without re-collecting audio, provided recording identifiers are recoverable and
each class retains enough independent recordings, the condition our capacity calculation checks.

\subsection{What is new here}

To our knowledge, score inflation under slice-level splitting and the feasibility of a
recording-level split have not been jointly quantified on a multi-library Chinese MIR benchmark.
Our three main arms control training volume while varying content origin; the CTIS-only C1
control probes non-test replacement. These controls narrow the interpretation, but do not fully
separate content identity from recording diversity.
